\section{Discussion（考察）}

\subsection{観光の共同生成化}

本研究で観察された観光体験は、「施設を利用する行為」ではなく、「他者と時間を共有する行為」としての側面を持っていた。

本事例では、観光客は「見る人」ではなく、「共に空間を作る人」となっていた。

特に「共に歌う」という行為は、

\begin{itemize}
\item 空間
\item 感情
\item 時間
\end{itemize}

を共有する行為であり、観光体験をより深いものへ変化させていた。


\subsection{地域素材の再接続}

本研究で重要なのは、新しいものを大量に生み出したわけではない点である。

既にそこには、

\begin{itemize}
\item 教会空間
\item オルガン
\item 音楽
\item 高齢者
\item 観光客
\item ホテル
\item クリスマス文化
\end{itemize}

が存在していた。

本企画で行われたのは、それらを適切に接続し直したことである。

これは地域創生において、「新規性」よりも「再接続」が重要である可能性を示している。


\subsection{情報設計としての地域接続}

本研究では、歌詞カード・曲解説・感謝メッセージなどを、人と文化と空間を接続する「情報設計」として扱った。

ここでいう情報設計とは、単なる情報伝達ではない。

\begin{itemize}
\item なぜ歌うのか
\item なぜ集まるのか
\item なぜ感謝するのか
\end{itemize}

という意味そのものを共有し、人間同士を自然に接続するための社会的設計を含んでいる。

本研究では、その情報設計によって、観光客・地域住民・ホテル関係者の境界が緩やかにつながっていく様子が観察された。


\subsection{高齢者の役割再形成}

本研究では、高齢者を単なる「支援対象」としてではなく、地域文化を支える中心的存在として捉えた。

高齢者が持つ経験や感情表現は、長年の地域生活によって蓄積された「地域知」とも捉えられる。

特に合唱における歌声は、人生経験そのものが音として現れる表現であり、その存在が空間全体へ深みを与えていた。

また、若い世代側が「その経験が必要である」と依頼することで、高齢者自身にも「まだ役割がある」という感覚が生まれ、社会参加への勇気や意味形成につながる可能性が示された。
